\documentclass[12pt,a4paper]{article}
\usepackage[utf8]{inputenc}
\usepackage[russian]{babel}
\usepackage{amsmath}
\usepackage{amsfonts}
\usepackage{listings}
\usepackage{hyperref}

\title{\textbf{ChainDB: Архитектурный и Научный Журнал Разработки}}
\author{Инженерный Совет Проекта ChainDB}
\date{Июль 2026}

\begin{document}

\maketitle
\tableofcontents
\newpage

\section{Введение и Философия Системы}
Проект \textbf{ChainDB} представляет собой исследование в области построения сверх-отзывчивых систем управления базами данных класса In-Memory. В отличие от традиционных подходов, разделяющих вычисления по потокам операционной системы (Threads), ChainDB эксплуатирует модель строго однопоточного ядра, управляемого через легковесные продолжения (\textit{Continuations}).

\section{Теоретико-Информационный Багаж}
\subsection{Изоморфизм Карри-Ховарда и Логика Пирса}
В вычислительной топологии ChainDB оператор управления контекстом \texttt{call/cc} рассматривается не просто как системный примитив, а как программное воплощение классического закона Пирса в теории типов:
\begin{equation}
((A \to B) \to A) \to A
\end{equation}
Данная типизация эквивалентна закону снятия двойного отрицания $\neg\neg A \to A$ в классической логике. Это позволяет рантайму базы данных осуществлять безопасные нелокальные переходы (``путешествия во времени'') к прошлым состояниям вычислительного стека, восстанавливая инварианты памяти без накладных расходов.

\subsection{Анатомия Памяти Chez Scheme: Spaghetti Stack}
Проблема монолитного непрерывного стека в Си-подобных архитектурах заключается в необходимости полного копирования байт (\texttt{memcpy}) при фиксации контекста, что уничтожает кэш процессора L1/L2.

Рантайм Chez Scheme решает это за счет архитектуры связанных сегментов стека в куче (\textit{Spaghetti Stack}). При вызове \texttt{call/cc}:
\begin{enumerate}
    \item Текущий сегмент стека изолируется и помечается как иммутабельный объект в куче ($O(1)$).
    \item Указатель стека процессора (\texttt{RSP/ESP}) перенаправляется на новый пустой фрейм.
    \item Сборщик мусора (Garbage Collector) берет изолированный сегмент под автоматический контроль.
\end{enumerate}

\section{Реализованные Архитектурные Решения}
\subsection{Компонент: Живые Курсоры (Zero-Copy Iterators)}
Для предотвращения деградации производительности ядра при выполнении тяжелых операций сканирования индекса разработан механизм изолированных фабрик курсоров на базе лексических замыканий.

Каждый созданный курсор инкапсулирует внутренние переменные состояния \texttt{resume-cont} и \texttt{return-cont} внутри локального контейнера \texttt{letrec*}. Это полностью исключает взаимное повреждение памяти (\textit{Memory Corruption}) при конкурентной обработке запросов от разных сетевых клиентов.

\section{Дорожная Карта Развития}
\begin{itemize}
    \item \textbf{Этап 1}: Пошаговая интеграция Диспетчера задач с изолированными курсорами.
    \item \textbf{Этап 2}: Интеграция системного сетевого слоя ввода-вывода через неблокирующий механизм \texttt{register-input-port}.
    \item \textbf{Этап 3}: Внедрение пула рабочих потоков (\textit{Worker Threads}) для обслуживания дисковой персистентности и тяжелых внешних вычислений.
\end{itemize}

\end{document}

